% ---
% id: optics-2026-001
% title: Очень простая задача №1
% topic: observations
% subtopic: telescope-optics
% difficulty: 1
% source: camp
% event: Оптика I
% year: 2026
% grade: 9-10
% lang: ru
% ---
Телескоп-рефрактор системы Кеплера имеет относительное отверстие $f/5$
при фокусном расстоянии объектива $F = \qty{100}{см}$.
\begin{enumerate}
  \item Вычислите увеличение телескопа для окуляра с $f = \qty{1}{см}$.
        Чему равна «длина трубы» телескопа — оптическое расстояние между
        объективом и окуляром?
  \item Пусть в фокальной плоскости объектива находится ПЗС-матрица
        с размером пиксела $\qty{10}{мкм}$. Определите расстояние между
        центрами изображений двух звёзд на полученном снимке (в пикселах),
        если на небе они разделены угловым расстоянием $\asec{20}$.
\end{enumerate}
